\chapter[Indifference of Measure]{The Indifference\\of Measure}

\vspace{1em}
\hfill
\begin{minipage}[c]{0.55\textwidth}
\raggedleft
\small\itshape
I am large, I contain multitudes.\\[0.5em]
\upshape\footnotesize Walt Whitman, ``Song of Myself'' (1855)
\end{minipage}%
\hspace{0.5em}
\begin{minipage}[c]{0.12\textwidth}
% \includegraphics[width=\textwidth]{figures/ch09_indifference.png}
\end{minipage}
\vspace{1.5em}

\noindent Part I built the framework. The thesis is what falls out of it.

The Library contains every reality. The universal prior is a measure on it. Your data restricts you to the subset of realities consistent with what you have observed. Within that subset, you exist at many indices. The measure has no preference about which one is you.

This is the book's title.

\section*{What the indifference is}

It is structural. The measure assigns probability weight to observer-moments. It does so by description length, computational depth, or whatever combination one specifies. None of these weights singles out a particular observer-moment as ``really you'' in any privileged sense.

Within the subset of realities consistent with your data, there are many candidates for \emph{you}. Some are in universes with simple laws; some in universes with complicated ones. Some are embedded observers with long lawful histories; some are momentary fluctuations. The measure has views about how to assign weight across these candidates. It has no view about which is yours.

You are a label drawn around a region of the measure. The measure draws no label.

\section*{The deeper indifference}

A parallel result is already familiar from physics. Special relativity tells us that ``now'' is not a fact about the universe but a fact about the observer: two observers moving relative to each other slice spacetime into different presents, and neither slicing is privileged. The natural reading of this is the \emph{block universe}: spacetime is a four-dimensional structure in which past, present, and future all exist at their coordinates. Your birth and your death are equally real. The structure contains no preferred temporal slice. This is the indifference of geometry. The argument is developed at length in the companion volume, \textit{Worldlines}; this book does not assume the reader has met it there.\footnote{The two books form a pair. \textit{Worldlines} works through the indifference of geometry from special relativity; this book works through the indifference of measure from algorithmic information theory and quantum mechanics. Either can be read alone. Read together, they make the same kind of argument at two levels of structure.}

This book establishes a deeper indifference. The block universe was indifferent to which moment in your worldline (your trajectory through spacetime) is ``now.'' The Library is indifferent to which observer-moment in your class of patterns is ``you.'' Same kind of indifference, one level further out.

The block universe dissolved the unique present. The Library dissolves the unique self.

\section*{What this is and is not}

The indifference is not an erasure. Your experience now is what it is; the moment of reading these words contains what it contains; the structure of your situation is real at the coordinates where you sit. The multiplicity does not undo any of it.

Nor is it annihilation. The pattern that constitutes ``you'' is instantiated, and each instantiation is real at the coordinates where it occurs. Multiplication does not subtract. If anything, it does the opposite: the structure is at many coordinates, each real, each containing whatever it contains.

The label \emph{I} also keeps its meaning. It picks out a region of the measure with certain features in common: a memory structure, a sensory state, a continuity of attention. It does not pick out a metaphysically unique entity, but it does not need to.

Nor does it answer \emph{where am I?} That question remains open. The framework permits the question and gives no sharp answer; the literature's attempts (SSA, SIA, the speed prior) all have costs the book chose not to pay. The question is left where it lives, unresolved.

The indifference, then, is a structural fact about the Library and the measure on it. The measure assigns weight to observer-moments. The weights have no opinion about which moment is yours.

\section*{What this does not console}

None of what follows in this book is consolation.

The framework is what it is. The measure is indifferent. The hellscapes from Chapter 8 are real at their coordinates; the Boltzmann brain hypothesis is open; the framework permits sustained suffering at indices we cannot reach, observers being deceived at coordinates the cut principle forbids us to exclude, versions of you and everyone you love living through anything physically realizable.

The remaining chapters do not argue these away. They walk into what the framework permits and ask: given this, what does it look like to be a finite indexed observer, with a life and people one loves, in a Library that does not single any of it out?

The response that comes is not derived from the mathematics. The mathematics leaves room. The room is what we walk into.

And the mathematics itself might be wrong. The framework takes three interpretive moves; any of them could fail. The response we choose does not depend on which way the cosmology turns out. The room is what we walk into either way.

\section*{Where this leaves you}

The thesis is stated. The book's title is on the page.

The next five chapters trace what the indifference does to continuation, death, love, the self, and what to do anyway. The arguments will follow from the framework already in place. The reader who has followed the cascade has all the tools.

The next chapter takes up continuation.
